\looseness=-1 Intelligent embodied agents need to quickly adapt to new scenarios by integrating long histories of experience into decision-making. For instance, a robot in an unfamiliar house initially wouldn't know the locations of objects needed for tasks and might perform inefficiently. However, as it gathers more experience, it should learn the layout of its environment and remember where objects are, allowing it to complete new tasks more efficiently. To enable such rapid adaptation to new tasks, we present \method, a new approach for in-context reinforcement learning (RL) for embodied agents. With \method, agents are capable of adapting to new environments using 64,000 steps of in-context experience with full attention while being trained through self-generated experience via RL. We achieve this by proposing a novel policy update scheme for on-policy RL called ``partial updates'' as well as a Sink-KV mechanism that enables effective utilization of a long observation history for embodied agents. Our method outperforms a variety of meta-RL baselines in adapting to unseen houses in an embodied multi-object navigation task. In addition, we find that \method is capable of few-shot imitation learning despite never being trained with expert demonstrations. We also provide a comprehensive analysis of \method, highlighting that the combination of large-scale RL training, the proposed partial updates scheme, and the Sink-KV are essential for effective in-context learning.
The code for \method and all our experiments is at \href{https://github.com/aielawady/relic}{github.com/aielawady/relic}.
